\section*{Chapitre \ref{ch:g12:kinematics-2d} --- Cinématique à deux dimensions}

\begin{solution}{exo:g12:kinematics-2d:1}
$\vect v = (3.0, 4.0)$, constant~; \omterm{def:g12:kinematics-2d:velocity}{vitesse} $\sqrt{3.0^2 + 4.0^2} =
\qty{5.0}{m/s}$. \omterm{def:g10:relative-motion:trajectory}{Trajectoire}~: la droite $y = \tfrac43 x$. Vecteur
\omterm{def:g12:kinematics-2d:velocity}{vitesse} constant~: \omterm{def:g12:kinematics-2d:uniform}{mouvement rectiligne uniforme}.
\end{solution}

\begin{solution}{exo:g12:kinematics-2d:2}
Droite, espacements égaux~: \omterm{def:g12:kinematics-2d:uniform}{rectiligne uniforme}~;
$v = 0.35/0.10 = \qty{3.5}{m/s}$. Des espacements qui s'ouvrent
signifieraient que la \omterm{def:g12:kinematics-2d:velocity}{vitesse} augmente.
\end{solution}

\begin{solution}{exo:g12:kinematics-2d:3}
$v(t) = 4.0\,t$, $a = \qty{4.0}{m/s^2}$~: accélération constante,
rectiligne uniformément accéléré. $v(3.0) = \qty{12}{m/s}$.
\end{solution}

\begin{solution}{exo:g12:kinematics-2d:4}
$f = 1/8.0 = \qty{0.125}{Hz}$~;
$v = 2\pi \times 4.5/8.0 \approx \qty{3.5}{m/s}$~;
$a = v^2/R = 3.53^2/4.5 \approx \qty{2.8}{m/s^2}$, pointant vers
l'axe du manège.
\end{solution}

\begin{solution}{exo:g12:kinematics-2d:5}
Pente~: $a = (10.0 - 2.0)/4.0 = \qty{2.0}{m/s^2}$. Aire (trapèze)~:
$d = \tfrac12 (2.0 + 10.0) \times 4.0 = \qty{24}{m}$.
\end{solution}

\begin{solution}{exo:g12:kinematics-2d:6}
$v = \qty{88.9}{m/s}$~: $a_n = 88.9^2/6000 \approx \qty{1.3}{m/s^2}
\approx 0.13\,g$. Comme $a_n = v^2/R$, une $v$ élevée exige un $R$
énorme pour garder l'accélération latérale (et la \omterm{def:g10:inertia:force}{force} sur la voie
et les passagers) acceptable --- le dévers change qui la fournit, non
sa taille.
\end{solution}

\begin{solution}{exo:g12:kinematics-2d:7}
$v(t) = -2.4\,t + 8.0$~; $a = \qty{-2.4}{m/s^2}$, constante.
Retournement à $v = 0$~: $t = 8.0/2.4 \approx \qty{3.3}{s}$, en
$x \approx \qty{16.3}{m}$. Avant~: en avant, freinant~; après~: en
arrière, accélérant --- uniformément accéléré tout du long.
\end{solution}

\begin{solution}{exo:g12:kinematics-2d:8}
$v_1 = 0.28/0.40 = \qty{0.70}{m/s}$~;
$v_2 = (0.54 - 0.10)/0.40 = \qty{1.10}{m/s}$~;
$v_3 = (0.88 - 0.28)/0.40 = \qty{1.50}{m/s}$. La vitesse gagne
\qty{0.40}{m/s} toutes les \qty{0.20}{s}~: $a = \qty{2.0}{m/s^2}$,
constante --- rectiligne uniformément accéléré.
\end{solution}

\begin{solution}{exo:g12:kinematics-2d:9}
$T = \qty{2.36e6}{s}$~: $v = 2\pi \times \num{3.84e8}/\num{2.36e6}
\approx \qty{1.0e3}{m/s}$~;
$a_n = v^2/R \approx \qty{2.7e-3}{m/s^2}$. Rapport~:
$g/a_n \approx 3600 = 60^2$ --- la gravité, diluée comme l'inverse du
carré de la distance (\cref{ch:g10:universal-gravitation}), rend
exactement compte du virage de la Lune.
\end{solution}

\begin{solution}{exo:g12:kinematics-2d:10}
$f = 1200/60 = \qty{20}{Hz}$, $T = \qty{0.050}{s}$~;
$v = 2\pi R f = 2\pi \times 0.25 \times 20 \approx \qty{31}{m/s}$~;
$a_n = v^2/R \approx \qty{3.9e3}{m/s^2} \approx 400\,g$. La paroi du
tambour fournit l'$a_n$ du linge~; à travers les trous rien ne pousse
l'eau vers l'intérieur, donc elle s'envole le long de la tangente.
\end{solution}

\begin{solution}{exo:g12:kinematics-2d:11}
$v_0 = \qty{25}{m/s}$~: $t_s = 25/7.5 \approx \qty{3.3}{s}$~;
$d = v_0^2/(2a) = 625/15 \approx \qty{42}{m}$. Graphique~: un segment
droit de \qty{25}{m/s} jusqu'à $0$ en $t_s$.
\end{solution}

\begin{solution}{exo:g12:kinematics-2d:12}
$\vect v = (-R\omega\sin\omega t,\, R\omega\cos\omega t)$, de
norme constante $R\omega$~; son produit scalaire avec $\vect{OM}$
s'annule, donc il est perpendiculaire au rayon, d'où tangent. En
dérivant encore, $\vect a = -\omega^2\vect{OM}$~: vers $O$, de norme
$\omega^2 R = (v/R)^2 R = v^2/R$.
\end{solution}

\begin{solution}{exo:g12:kinematics-2d:13}
$a_t = \qty{2.5}{m/s^2}$ (vers l'arrière), $a_n = 20^2/80 =
\qty{5.0}{m/s^2}$~;
$a = \sqrt{2.5^2 + 5.0^2} \approx \qty{5.6}{m/s^2}$~; angle avec la
direction du mouvement~: $\theta = 180^\circ - \arctan(5.0/2.5) \approx 117^\circ$
(c.-à-d.\ $63^\circ$ depuis l'arrière droit). $5.6 < 7.0$~: dans
l'adhérence, avec peu de marge.
\end{solution}

\begin{solution}{exo:g12:kinematics-2d:14}
Lancement~: $t_1 = 20/1.2 \approx \qty{17}{s}$,
$d_1 = 20^2/(2 \times 1.2) \approx \qty{167}{m}$. Croisière~:
\qty{40}{s}, $\qty{800}{m}$. Freinage~: $t_3 = 20/1.5 \approx
\qty{13}{s}$, $d_3 = 400/3.0 \approx \qty{133}{m}$. Total~:
$\qty{1.1e3}{m}$ en \qty{70}{s}~; moyenne $1100/70 \approx
\qty{16}{m/s}$ (\qty{57}{km/h}). $v(t)$ est un trapèze dont l'aire,
$\tfrac12(70 + 40) \times 20 = \qty{1100}{m}$, vérifie le total.
\end{solution}

\begin{solution}{exo:g12:kinematics-2d:15}
$v = 2\pi \times \num{4.22e7}/\num{86164} \approx \qty{3.1e3}{m/s}$~;
$a_n = v^2/R \approx \qty{0.22}{m/s^2}$. La gravité terrestre à cette
distance la fournit~; la \omterm{def:g12:kinematics-2d:ucm}{période} égale la
\omterm{def:g12:kinematics-2d:ucm}{période} de rotation de la Terre, donc
\omterm{def:g10:universal-gravitation:satellite}{satellite} et sol tournent ensemble et il reste au-dessus
d'un point équatorial (\cref{ch:g12:satellites-kepler}).
\end{solution}

\begin{solution}{pb:g12:kinematics-2d:1}
\textbf{1.} Le mouvement est relatif (\cref{ch:g10:relative-motion})~:
chaque position et vitesse doit nommer son référentiel. Ici~: le
référentiel du sol (route), $x$ le long de la route depuis la position
de la première image.

\textbf{2.} Chaque intervalle~: $9.0/0.50 = \qty{18.0}{m/s}$ --- quatre
valeurs égales.

\textbf{3.} Route droite, déplacements égaux en temps égaux~:
\omterm{def:g12:kinematics-2d:uniform}{mouvement rectiligne uniforme}.

\textbf{4.} $v_0 = \qty{18.0}{m/s} = \qty{64.8}{km/h}$.

\textbf{5.} $\vect a = \vect 0$.

\textbf{6.} $v(t) = -a\,t + v_B$~; $x(t) = -\tfrac12 a t^2 + v_B\,t$.

\textbf{7.} $v(t_s) = 0$ donne $t_s = v_B/a$~; alors
$d_b = -\tfrac12 a (v_B/a)^2 + v_B^2/a = v_B^2/(2a)$.

\textbf{8.} $v_B = \sqrt{2 a d_b} = \sqrt{2 \times 6.0 \times 27} =
\sqrt{324} = \qty{18.0}{m/s}$ --- exactement la
\omterm{def:g12:kinematics-2d:velocity}{vitesse} filmée~: le conducteur n'a jamais ralenti avant
de freiner.

\textbf{9.} Réaction~: $18.0 \times 1.0 = \qty{18}{m}$~; total
$18 + 27 = \qty{45}{m}$.

\textbf{10.} $d_b = v_B^2/(2a) \propto v_B^2$~: doubler $v_B$ multiplie
$d_b$ par $4$. À \qty{50}{km/h} ($\qty{13.9}{m/s}$)~:
$d_b = 13.9^2/12.0 \approx \qty{16}{m}$.

\textbf{11.} La \omterm{def:g12:kinematics-2d:velocity}{vitesse} est constante mais la direction de la
vitesse tourne~: $\vect v$ n'est pas constant, donc la voiture
accélère.

\textbf{12.} Quart de tour~: $2\pi R/4 = \qty{19.6}{m}$~;
$v = 19.6/2.5 \approx \qty{7.9}{m/s}$ (\qty{28}{km/h}).

\textbf{13.} $a_n = 7.85^2/12.5 \approx \qty{4.9}{m/s^2}$, vers le
centre du rond-point.

\textbf{14.} $a_{\max} = \mu g = 0.70 \times 9.81 \approx
\qty{6.9}{m/s^2}$~; $v_{\max} = \sqrt{\mu g R} =
\sqrt{6.87 \times 12.5} \approx \qty{9.3}{m/s}$ (\qty{33}{km/h}).
$7.9 < 9.3$~: dans l'adhérence --- cohérent avec la voie propre, sans
dérapage.

\textbf{15.} $T = 2\pi R/v = 78.5/7.85 \approx \qty{10.0}{s}$~;
$f = \qty{0.10}{Hz}$.

\textbf{16.} $64.8 - 50 \approx \qty{15}{km/h}$ de trop --- environ
$30\%$ au-dessus de la limite.

\textbf{17.} Le freinage agit sur $41 - 18 = \qty{23}{m}$~:
$v^2 = v_B^2 - 2ad = 324 - 2 \times 6.0 \times 23 = 48$, donc
$v \approx \qty{6.9}{m/s} \approx \qty{25}{km/h}$ à l'impact.

\textbf{18.} À \qty{13.9}{m/s}~: $13.9 \times 1.0 + 13.9^2/12.0 =
13.9 + 16.1 = \qty{30.0}{m} < \qty{41}{m}$ --- la voiture s'arrête
\qty{11}{m} avant la fourgonnette.

\textbf{19.} $d(v) = v\,t_r + v^2/(2a)$~: une fois linéairement (la
réaction couvre $v t_r$ avant que les freins n'agissent), une fois au
carré (la partie cinétique $v^2/(2a)$) --- d'où le coût
disproportionné des excès modérés.

\textbf{20.} La dashcam et les traces s'accordent~: \qty{65}{km/h} dans
une zone à \qty{50}{km/h}~; à la limite légale la voiture se serait
arrêtée \qty{11}{m} avant la fourgonnette --- l'accident, c'est les
\qty{15}{km/h} de trop.
\end{solution}
